\section{Differentiable MDP via Occupancy Measure LP}
We now reformulate the MDP as a linear program over occupancy measures and derive a closed-form gradient for the DFL objective. This replaces the nonlinear Bellman equation with linear constraints, and the resulting gradient depends only on a small basis matrix.

\subsection{LP Formulation of Discounted MDPs}
The discounted state-action occupancy measure $y \in \R^{|\gS||\gA|}$ is defined as $y(s,a) = \sum_{t=0}^\infty \beta^t \Pr(s_t = s, a_t = a)$, the expected discounted frequency of visiting state $s$ and taking action $a$ over the infinite horizon. It is well-known that the MDPs can be equivalently written as the following \LP~\cite{}:
\begin{align}
    y^*_\theta = \arg\max_{y \geq 0} 
    r_\theta^\top y \quad
    \text{s.t.} \quad
    H_\theta\, y = \gamma, \tag{\texttt{LP-1}}
    \label{eq:lp_primal}
\end{align}
where $r_\theta \in \mathbb{R}^{|\mathcal{S}||\mathcal{A}|}$ is the reward
vector and
$H_\theta \in \mathbb{R}^{|\mathcal{S}| \times |\mathcal{S}||\mathcal{A}|}$
is the flow constraint matrix with entries
\begin{align}
    H_{s,(s',a')} = \delta_{s,s'} - \beta\, P_\theta(s \mid s',a'),
    \label{eq:H_def}
\end{align}
where $\delta_{s,s'}$ is the Kronecker delta, i.e., $\delta_{s,s'} = 1$ iff $s = s'$ otherwise $\delta_{s,s'} = 0$. The optimal policy can be recovered from $y^*_\theta$ by $\pi^*_\theta(a|s) = y^*_\theta(s, a)/\sum_{a' \in \gA} y^*_\theta(s, a')$.

\paragraph{Comparison with prior DFL gradients.}
Existing DFL approaches for MDPs differentiate through the KKT conditions of the Bellman equation and treat the optimal policy as the implicit solution of a \emph{nonlinear} fixed-point system~\citep{wang2021learning, amos2017optnet}. This requires forming and inverting a Hessian-like matrix of dimension $|\gS||\gA| \times |\gS||\gA|$, whose cost grows quickly with the action space and limits scalability beyond small MDPs. The \LP reformulation differs in two important ways. First, the optimality conditions become \emph{linear} in the decision variable $\bm{y}$, so the sensitivity of $\bm{y}^*_{\bm{\theta}}$ admits a closed-form expression rather than an implicit-function expansion. Second, under non-degeneracy only the basic component $\bm{y}^B_{\bm{\theta}} \in \R^{|\gS|}$ contributes to the gradient (\Cref{eq:dfl_gradient}), and thus the dominant linear solve has dimension $|\gS| \times |\gS|$, independent of $|\gA|$. As we show in \Cref{tbl:1}, this yields an order-of-magnitude reduction in compute on standard benchmarks, which in turn enables the soft state aggregation extension introduced in~\Cref{sec:agg}.
%---------- 
\subsection{Non-Degeneracy Property and the Active Basis}
In Linear Programming, a basic feasible solution is said to be \emph{degenerate} if one or more of its basic variables are zero. Degeneracy can cause several complications, since multiple bases may correspond to the same extreme point of the feasible polytope. However, for the linear program $\mathtt{LP}$ defined over occupancy measures, if the initial distribution is strictly positive everywhere ($\bm{\gamma}(s) > 0$ for all $s \in \mathcal{S}$), then any basic feasible solution is guaranteed to be \emph{non-degenerate}. This implies that the optimal solution $\bm{y}_{\bm{\theta}}^*$ contains exactly $|\mathcal{S}|$ non-zero entries (the basic variables). Equivalently, for each state $s \in \mathcal{S}$ there is exactly one action $a \in \mathcal{A}$ with strictly positive occupancy measure $y(s,a)$, corresponding to the unique optimal action $\pi_{\bm{\theta}}^*(s)$ in the optimal policy. As a consequence, the associated basis matrix, denoted by $\bm{H}_{\bm{\theta}}^B$, is unique, square, and invertible with dimension $|\mathcal{S}| \times |\mathcal{S}|$. This invertibility is the mathematical foundation that ensures the sensitivity derivative is well-defined and numerically stable. In contrast, under degeneracy, multiple different basis matrices may correspond to the same solution, making differentiation problematic. 


%----------
\paragraph{Fixed-basis Assumption.}
We assume that for a small perturbation in the parameters $\bm{\theta}$, the optimal policy $\pi_{\bm{\theta}}^*$ does not change. This implies that the set of non-zero indices in $\bm{y}_{\bm{\theta}}^*$ remains constant, and consequently, the optimal basis matrix $\bm{H}_{\bm{\theta}}^B$ also remains fixed locally. Under this assumption, the mapping from $\bm{\theta}$ to $\bm{y}_{\bm{\theta}}^*$ is linear and differentiable, allowing us to compute gradients without encountering discrete jumps in the policy.

\paragraph{Practical considerations.}
In practice, the optimal basis may switch as $\bm{\theta}$ is updated during training, especially in early epochs when the predicted parameters are far from the ground truth. At a switching point the gradient~\eqref{eq:dfl_gradient} is not formally defined, since two distinct bases attain the LP optimum. We handle this in two complementary ways. First, the augmented Lagrangian penalty introduced in~\Cref{eq:dfl_loss} below provides a smooth surrogate signal that remains informative across basis transitions. Second, we initialize the predictor with a brief prediction-loss warm-up so that joint training begins in a region where the basis is comparatively stable. We further track the empirical frequency of basis changes per epoch as a diagnostic, and find that, after the warm-up phase, basis switches become rare for the MDP instances we consider. The closed-form gradient is therefore an accurate approximation of the true sensitivity in the regime where the predictor has been adequately initialized.

%--------------------
\subsection{Closed-form Gradient of the \LP Layer}
%--------------------

The \LP objective is $v^*_\theta = r_\theta^\top y^*_\theta$. 
Using the product rule, the total derivative is the sum of the direct gradient (change in rewards) and the indirect gradient (change in occupancy), as follows:
\begin{align}
    \frac{d v_{\theta}^*}{d\bm{\theta}}
    & =
    \underbrace{\left( \frac{\partial \bm{r}_{\bm{\theta}}}{\partial \bm{\theta}} \right)^\top 
    \bm{y}_{\bm{\theta}}^*}_{\text{Direct Term}}
    +
    \underbrace{\bm{r}_{\bm{\theta}}^\top 
    \left( \frac{d \bm{y}_{\bm{\theta}}^*}{d \bm{\theta}} \right)}_{\text{Indirect Term}}.
    \label{eq:lp_grad}
\end{align}

The direct term can be easily derived. In order to compute the indirect term, we exploit the 
relationship between $\bm{y}_{\bm{\theta}}^*$ and $\bm{\theta}$ defined by the linear constraints 
that prescribe the polyhedral domain of the space of occupational measures,
$\bm{H}_{\bm{\theta}} \bm{y}_{\bm{\theta}}^* = \bm{\gamma}$. 
We assume that the initial distribution $\bm{\gamma}$ is independent of $\bm{\theta}$, that is, 
$\frac{d\bm{\gamma}}{d\bm{\theta}} = \bm{0}$, and therefore
\begin{align*}
    \frac{\partial \bm{H}_{\bm{\theta}}}{\partial \bm{\theta}} \bm{y}_{\bm{\theta}}^*
    + \bm{H}_{\bm{\theta}} \frac{d \bm{y}_{\bm{\theta}}^*}{d \bm{\theta}}
    = \bm{0}.
\end{align*}
Partitioning $y^*_\theta = (y^B_\theta, y^N_\theta)$ and $H^*_\theta = (H^B_\theta, H^N_\theta)$ 
into basic and non-basic components, non-degeneracy implies that $y^B_\theta \in \R^{|\gS|}$, $y^N_\theta = 0$ with $H^B_\theta \in \R^{|\gS| \times |\gS|}$ invertible.
Under the fixed basis assumption, we establish
\begin{align}
    \frac{d \bm{y}_{\bm{\theta}}^B}{d \bm{\theta}}
    & = - \left(\bm{H}_{\bm{\theta}}^B\right)^{-1}
    \left( \frac{\partial \bm{H}_{\bm{\theta}}^B}{\partial \bm{\theta}} \bm{y}_{\bm{\theta}}^B \right).
    \label{eq:jaco}
\end{align}

Substituting \Cref{eq:jaco} into \Cref{eq:lp_grad} and noting that
$\left( \frac{\partial \bm{r}_{\bm{\theta}}}{\partial \bm{\theta}} \right)^\top 
    \bm{y}_{\bm{\theta}}^* 
    = \left( \frac{\partial \bm{r}_{\bm{\theta}}^{B}}{\partial \bm{\theta}} \right)^\top 
    \bm{y}_{\bm{\theta}}^B$,
we obtain
\begin{align}
    \frac{d v_{\bm{\theta}}^*}{d\bm{\theta}}
    &=
    \left( \frac{\partial \bm{r}_{\bm{\theta}}^{B}}{\partial \bm{\theta}} \right)^\top 
    \bm{y}_{\bm{\theta}}^B
    -
    \left(\bm{r}_{\bm{\theta}}^{B}\right)^\top
    \left( \bm{H}_{\bm{\theta}}^{B} \right)^{-1}
    \left( 
        \frac{\partial \bm{H}_{\bm{\theta}}^B}{\partial \bm{\theta}} 
        \, \bm{y}_{\bm{\theta}}^B
    \right).
    \label{eq:dfl_gradient}
\end{align}
This formula provides the exact gradient of the expected reward with respect to the MDP parameters for the current optimal policy.

\paragraph{Training loss.}
Rather than minimizing the regret in~\Cref{eq:dfl_obj} directly, we adopt an augmented Lagrangian surrogate that combines a decision-quality term with a soft feasibility penalty:
\begin{align}
    L(\bm{y}^*_{\bm{\theta}})
    := -\,\bm{r}_{\mathrm{true}}^\top \bm{y}^*_{\bm{\theta}}
       + \rho\, \big\|\bm{H}_{\mathrm{true}}\, \bm{y}^*_{\bm{\theta}} - \bm{\gamma}\big\|^2.
    \label{eq:dfl_loss}
\end{align}
The first term is the negative downstream value evaluated at the prediction-induced occupancy $\bm{y}^*_{\bm{\theta}}$ under the \emph{ground-truth} reward and transition matrix, and minimizing it directly drives the predictor to maximize true decision quality. The second term penalizes violation of the true flow constraints, which in general are not satisfied because $\bm{y}^*_{\bm{\theta}}$ is feasible only with respect to the predicted dynamics $\bm{H}_{\bm{\theta}}$. This penalty serves two complementary roles. First, it provides a continuous signal across changes of the optimal basis discussed above, and thereby stabilizes training in the regime where the fixed-basis assumption is locally violated. Second, it injects useful curvature into the loss landscape even when the LP value is locally flat in $\bm{\theta}$, which is a known failure mode of differentiable optimization layers. The coefficient $\rho > 0$ balances decision quality and constraint adherence; we use $\rho = 0.1$ throughout our experiments and find that the result is insensitive to $\rho$ within an order of magnitude. Optionally, a small MSE prediction loss $\lambda \cdot \|m_w(x) - \theta^*\|^2$ can be added during early epochs to anchor the predictor to the ground truth, providing a curriculum from prediction to decision.

%--------------------
\subsection{Computational Complexity Comparison}
%--------------------
Table \ref{tbl:1} compares the computational cost of the LP/Occupancy Measure approach derived above with the KKT/Policy Gradient approach (as used in Wang et al., 2022), along with an example for the specific case of a \textbf{Gridworld} problem on a $5 \times 5$ grid with 25 states and 5 actions. The dominant cost in both methods is matrix inversion.
\begin{table}[ht!]
    \centering
    \caption{\small{Comparison of computational complexity between the LP  approach and the KKT approach.}} 
    \vspace*{0.5em}
    \label{tbl:1}
    
    \begin{tabular}{l|c|c}
        \textbf{Metric} &
        \textbf{LP Approach} &
        \textbf{KKT Approach} \\[0.1em] \hline 
        Core Matrix to Invert &
        Optimal Basis Matrix $\bm{H}_{\bm{\theta}}^B$ &
        Hessian Matrix $\nabla_{\pi}^2 J(\pi)$ \\[0.1em]
        Matrix Dimension &
        $|\mathcal{S}| \times |\mathcal{S}|$ &
        $|\mathcal{S}||\mathcal{A}| \times |\mathcal{S}||\mathcal{A}|$ \\[0.1em]
        Gridworld Dimensions &
        $25 \times 25$ &
        $125 \times 125$ \\[0.1em]
        Inversion Complexity &
        $\mathcal{O}(25^{2.37})$ &
        $\mathcal{O}(125^{2.37})$ \\[0.1em]
        Approximate Operations &
        $2{,}056$ &
        $93{,}256$
    \end{tabular}
\end{table}

\begin{remark}
    When a pivoting algorithm such as the Simplex method is used to solve $\mathtt{LP}$, the inverse of the final basis matrix is directly available from the last tableau at no additional computational cost. Moreover, although the Simplex algorithm has exponential worst--case complexity for general linear programs, more recent results show that, when solving $\mathtt{LP}$ with a fixed discount factor $\beta$, the Simplex method in fact admits a polynomial time complexity bound. This observation can substantially reduce the computational burden when incorporating $\mathtt{LP}$ into a DFL framework.
\end{remark}

\section{State Aggregation}
\label{sec:agg}
The \LP gradient scales with $|S|$, which remains the bottleneck when the state space is large or continuous. We address this with a soft state-aggregation layer that is trained jointly with the predictor and extend to the general state MDPs through linear function approximation of occupancy measures.

\subsection{Soft State Aggregation for Finite State MDPs}
To make the LP tractable for large state spaces, we introduce a \emph{soft state aggregation} scheme that compresses the original $S$-state MDP into a reduced MDP with $M \ll S$ abstract states. Unlike hard clustering approaches that assign each state to a single cluster, we maintain a differentiable \emph{membership matrix} $W \in \R^{S \times M}$, where each row lies on the probability simplex ($W \geq 0$, $\sum_m W_{sm} = 1$ for all $s$). The entry $W_{sm}$ represents the soft probability that ground state $s$ belongs to abstract state $m$. This soft assignment is crucial: it enables end-to-end gradient-based optimization of the aggregation structure jointly with the downstream decision objective, whereas hard clustering would produce a piecewise-constant loss surface with respect to the aggregation parameters.

\paragraph{Learnable membership construction.}
We parameterize the membership matrix through aggregation parameters $\xi = (E, \Theta_W)$, where $E \in \R^{S \times d_e}$ is a state embedding matrix and $\Theta_W$ maps features to cluster logits. Specifically, we first compute a transition-aware feature for each state-action pair as the expected next-state embedding $\phi(s, a) = \sum_{s'} P(s' \mid s, a) E_{s'} \in \R^{d_e}$, then form a state-level representation by concatenating across all actions: $\bar\phi(s) = [\phi(s,1); \ldots; \phi(s,A)] \in \R^{A d_e}$. This feature encodes how each state interacts with the transition structure under every action. The soft cluster assignment is then obtained by projecting $\bar\phi(s)$ onto $M$ cluster logits via $\Theta_W$ and applying a temperature-scaled softmax:
\begin{align*}
    W_{sm}(\xi) = \frac{\exp\bigl((\Theta_W^\top \bar\phi(s))_m / \tau\bigr)}{\sum_{m'} \exp\bigl((\Theta_W^\top \bar\phi(s))_{m'} / \tau\bigr)}.
\end{align*}
Given the soft membership $W$, we construct a reduced MDP over $M$ abstract states by computing weighted averages of the original MDP quantities. Let $\mu_m = \sum_s W_{sm}$ denote the effective mass of cluster $m$. The aggregated rewards, transitions, and initial distribution are:
\begin{align*}
    \hat r_{ma} = \frac{1}{\mu_m} \sum_s W_{sm}\, r_{sa}, \quad
    \hat P_{man'} = \frac{1}{\mu_m} \sum_{s, s'} W_{sm} P(s' \mid s, a)\, W_{s'n'}, \quad 
    \hat\gamma_m = \sum_s \gamma_s\, W_{sm}.
\end{align*}
The aggregated flow constraint matrix $\hat H \in \R^{M \times MA}$ takes the same form as the original~\Cref{eq:H_def}:
\begin{align*}
    \hat{H}_{s,(s',a')} & = \delta_{s,s'} - \beta \hat{P}_{\theta}(s \mid s', a').
\end{align*}
We then solve the reduced LP with only $M$ state and $M|\gA|$ variables on the aggregated MDP:
\begin{align}
    \max_{\hat y \geq 0} \quad & \hat r_\theta^\top \hat y \quad
    \text{s.t.} \quad \hat H(\theta, \xi) \hat y = \hat\gamma(\xi). \tag{\texttt{LP-2}} \label{eq:agg_lp}
\end{align}
To recover state-action occupancy measures in the original state space, we lift the aggregated solution back via the membership matrix:
\begin{align*}
    y_{\text{disagg}}(\theta, \xi) = W\hat{y} (\theta, \xi).
\end{align*}
Since $W$ is differentiable in $\xi$, the aggregated form can be trained jointly with the prediction parameter $w$ by backpropagating through~\Cref{eq:dfl_gradient}

\subsection{Function Approximation for General State MDPs}
Soft state aggregation can be viewed as a projection of the occupancy \LP onto the column span of $W$. This perspective generalizes immediately to general state MDPs by replacing the clustering matrix with feature mappings. Specifically, 

Let the state space $\gS$ be continuous with initial distribution $\gamma$ and transition density $P(s | s', a')$. The discounted occupancy $y : \gS \times \gA \to \mathbb{R}_{+}$ satisfies the infinite-dimensional LP
\begin{align*}
    \max_{y \geq 0} \int_{\gS \times \gA} r_\theta(s, a) y(s, a) dsda
    \quad\text{s.t.}\quad
    (B - \beta P)y = \gamma, 
\end{align*}
where we define operators $(B y)(s) := \int_{\mathcal{A}} y(s, a) da$ and $(P y)(s) := \int p(s \mid s', a') y(s', a') ds' da'$. We approximate the occupancy $y$ by a finite expansion in $d$ state-action basis functions $\Phi := [\phi_1, \dots, \phi_d]: \gS \times \gA \to \R^d$, and project the flow constraint, which lives in $\gS$, onto $m$ state basis function $\Psi := [\psi_1, \dots, \psi_m] : \gS \to \R^m$,
\begin{align*}
    y_w(s,a) = \Phi w =  \sum_{j=1}^d w_j \phi_j(s,a), \quad \Psi^\top ((B - \gamma P)\Phi w - \rho) = 0.
\end{align*}
Note that we can write the above \emph{infinite number} of flow constraints into a linear system $Aw = b$ with
$A := \Psi^\top (B - \gamma P)\Phi \in \mathbb{R}^{m \times d},
b := \Psi^\top \rho \in \mathbb{R}^{m}$.
If the basis functions $\phi_j$ are non-negative, the constraint $w \geq 0$ enforces $y_w \geq 0$. The approximated LP is
\begin{align}
    \max_{w \geq 0}  \tilde{r}_\theta^\top w
    \quad\text{s.t.}\quad
    A w = b, .\tag{\texttt{LP-3}}
    \label{eq:projected_lp}
\end{align}
where $(\tilde{r}_\theta)_j = \int_{\mathcal{S}\times\mathcal{A}} r_\theta(s, a) \phi_j(s, a) ds da$. In practice, $\Phi$ and $\Psi$ can be realized by neural-network features or other parametric families such as random Fourier features~\cite{}, and they can be trained jointly with the predictor through~\Cref{eq:dfl_gradient}.

Here, we provide how to train the neural network features by using trajectories. Given trajectory samples $(s_t, a_t, s_{t+1})$, we can estimate operators in $A$ by
\begin{align*}
    (\Psi^\top P \Phi)_{ij} 
    \approx \frac{1}{N}\sum_{t=1}^N \psi_i(s_{t+1})\phi_j(s_t, a_t), \quad
    (\Psi^\top B \Phi)_{ij} \approx \frac{1}{N_s}\sum_t \psi_i(s_t)\sum_a \phi_j(s_t, a).
\end{align*}


\remark[Recovery of soft aggregation]
{If we choose $\phi_j(s, a) = W_{sm} \mathbf{1}[a = a']$ for a cluster $m$ and action $a'$, and $\psi_i(s) = W_{s i}$, then the projected LP~\eqref{eq:projected_lp} coincides with the finite aggregated LP~\eqref{eq:agg_lp}: $M = \hat H$, $b = \hat\gamma$, and $c = \hat r$. The continuous formulation thus strictly generalizes soft aggregation.}

